We specified Invariance-Based Architectural Investigation Protocol v1.0: a reproducible, theory-independent methodology for evaluating candidate architectural invariants. The protocol standardizes (i) observation--theory separation, (ii) an end-to-end research pipeline, (iii) required artifacts for replication, (iv) measurement and classification rules, and (v) protocol versioning.

A key contribution is the explicit separation of the protocol (Paper~0, the scientific instrument) from any single execution of it (e.g., the Three-System Pilot as B1). This decoupling enables a coherent research program in which successive studies (B2, B3, \ldots) can be conducted by independent teams using the same instrument and compared on a common basis.

Paper~0 establishes the shared baseline so that future invariant claims can be compared, replicated, and aggregated across research groups. The next objective is to raise the protocol's operational rigor to the point where an external group can execute B2 or B3 using only this document and produce comparable results. Future work therefore prioritizes tightening the replication kit (artifact checklists, reporting templates, and acceptance criteria) alongside reference schemas for common architectural representations and reusable extractor/normalization toolchains for major ecosystems.
